\section{Discussion}
\label{sec:discussion}

The three tasks do not support a single verdict, and the differences between
them are more informative than any average. The first version of this benchmark
reported a uniform classical win. Three of its ingredients do not survive the
revised protocol: the truth sat inside the parametric basis, the learned gap
filler was evaluated under a normalisation it was not trained with, and the
downstream noise model was re-estimated on denoised outputs. Correcting them
reverses the gap-filling verdict, makes the denoising verdict conditional, and
leaves detection negative.

\paragraph{The denoising advantage is a property of the truth family.}
On the in-family truth the trajectory model removes $0.790$ of the injected
error and the RTS smoother $0.791$, while the U-Net removes $0.366 \pm 0.002$
and the model-free decompositions lie between $0.232$ and $0.349$. That ordering
is expected: the truth is a trend plus constant harmonics with mild modulation,
which is almost exactly the basis of the parametric estimator. Moving the truth
outside that basis changes the ranking. The trajectory model falls to $-0.106$
and the RTS smoother to $-0.104$: on a signal it cannot represent, the fit adds
error rather than removing it. The model-free decompositions barely move, and
the U-Net, trained on the matching out-of-family generator, does not degrade at
all: it scores $0.363 \pm 0.001$. Scored with per-window normalisation, it is
above every model-free decomposition in both families. On the synthetic grid it
reaches $0.389 \pm 0.001$ in-family and $0.381 \pm 0.001$ out-of-family, against
$0.320$ and $0.318$ for the moving median and $0.317$ and $0.312$ for the
wavelet. The defensible statement is therefore narrower than a strong
learned-denoiser claim: the learned denoiser is the best non-parametric choice
in both families and the only learned arm that is nearly family-invariant, but
it remains far below the parametric estimator when the truth lies in its basis.
The earlier reading, in which the learned denoiser won only on the synthetic
grid and lost to the wavelet on real stations, was an artefact of whole-series
normalisation at inference.

\paragraph{Gap filling improves but does not win on the full real set.}
The first version evaluated the network with whole-series normalisation while
training it on per-window statistics. Scoring the same ten nets per window
removes the deficit: at 30 days the per-window U-Net reaches $3.07 \pm 0.03$~mm
on long synthetic records and $2.74 \pm 0.03$~mm on the 30-series real subset,
against $3.34$ and $2.93$ for linear interpolation and $3.79$ and $3.35$ for the
trajectory model. At 5 days it is $0.23$~mm behind SSA on synthetic records and
$0.05$~mm behind on real ones, and at 120 days the trajectory model remains
better ($3.99$ against $5.09 \pm 0.08$). The whole-series arm is not a
competitive configuration at any length. The correction does not extend to the
full held-out data: across all 201 test component series the U-Net ties linear
interpolation on NGL ($2.964$ against $2.967$~mm) and falls $0.06$~mm behind on
JPL ($3.293$ against $3.233$), and it leads by region only in Asia and Europe
(Sec.~\ref{sec:results:external}). The learned filler is therefore competitive
at intermediate gaps, and the first version's failure was an inference defect
rather than an architectural limit, but the corrected result is parity on real
data, not a product-level win. The defect is also a warning for the field: a
normalisation mismatch between training and inference can look like a
model-capacity failure, and a single-seed, single-protocol evaluation cannot
tell the two apart.

The mechanism behind the win is visible in the signal-error arm. The trajectory
model's error against the missing true signal stays near $0.7$~mm at every gap
length, because it extrapolates a low-order trend and harmonics; the
whole-series network reaches $5.97$~mm at 30 days and $61.0$~mm at 365 days.
The per-window network inherits the interpolation prior learned over gaps of
5--120 days, which is why it is competitive in the middle of that range and why
its advantage should not be extrapolated to year-long gaps, for which no
per-window measurement was made.

\paragraph{Detection: a better ranking does not give a better operating point.}
On synthetic offsets the classical detector has the higher validation-selected
$F_1$ at every amplitude, from $0.073$ against $0.0005$ at 2~mm to $0.975$
against $0.725$ at 25~mm. The learned field nevertheless has the higher average
precision at 16 and 25~mm ($0.648$ and $0.669$ against $0.208$ and $0.007$). The
two metrics answer different questions: average precision integrates over all
thresholds, while the deployed detector uses one threshold chosen on held-out
stations. The learned field orders large events well but transfers its operating
point poorly. On real catalogue labels the gap widens: classical $F_1$ is
$0.457 \pm 0.027$, the twenty U-Net seeds average $0.232 \pm 0.107$, the ten TCN
seeds $0.273 \pm 0.096$, and four seeds across the two architectures collapse to
essentially no detections. Where a learned seed does fire, its matched detections
sit closer to the catalogued epoch (absolute median delay $16.5$--$52.1$ days for
the U-Net and $7.8$--$39.9$ days for the TCN, against $69.7$ for the classical
detector), but it finds fewer events. The catalogue is not ground truth either: its dates may refer to
the earthquake rather than to the observable break, which is why the real-label
arm uses a 180-day tolerance. The comparison is nonetheless the one a geodesist
would make, and it does not show a learned advantage.

\paragraph{What the velocity product inherits.}
The corrected downstream test removes the earlier warning rather than changing
its mechanism. With the noise model estimated on the raw series and held fixed,
the trajectory MLE is unbiased ($-0.0024$~mm\,yr$^{-1}$) with a velocity RMSE of
$0.0586$~mm\,yr$^{-1}$, and the classical denoisers stay within
$0.0588$--$0.0616$. The U-Net denoiser reaches $0.0605$ (two seeds, $0.0600$ and
$0.0611$) with a bias of $-0.0034$, and the U-Net imputer is indistinguishable
from the classical fillers at $0.0586$. Calibration is $0.649$ for the
reference, $0.618$--$0.647$ for the classical denoisers and $0.628$ for the
learned denoiser, with 1$\sigma$ coverage of $0.389$, $0.444$--$0.472$ and
$0.458$ respectively. The values of the first version ($0.2087$ and $0.197$)
came from whole-series normalisation, which mismatches the statistics the
network was trained under; scoring window by window removes them, and the
learned arm is neither better nor worse than the classical denoisers by more
than $0.002$~mm\,yr$^{-1}$. Re-estimating the noise on a denoised output, the
other protocol of the first version, still collapses the reported uncertainty
for the trajectory and RTS outputs because they leave essentially no residual;
that arm is reported to explain the earlier number, not as a valid pipeline.
Adding a random-walk component lowers the reference RMSE from $0.0586$ to
$0.0690$ and the calibration ratio from $0.649$ to $0.575$; the U-Net denoiser
reaches $0.0698$ and the imputer $0.0690$, and the ordering is unchanged. The
product-level rule is therefore weaker than in the first version: no method,
learned or classical, improves a velocity that the joint trajectory/noise fit
already estimates, but learned denoising no longer damages the product. The
decision should be made on the product metric, not on the pointwise signal
error, because the two do not move together.

\paragraph{Architecture and seeds.}
The detection result is not a statement about one architecture or one draw. On
the real labels the seed standard deviation is $0.107$ for the U-Net and
$0.096$ for the TCN, nearly half the mean in both cases, and on synthetic
offsets at 8~mm the U-Net seed distribution has a standard deviation of $0.222$
around a mean of $0.436$. A TCN does better than the U-Net at every amplitude
above 4~mm ($0.829$ against $0.686$ at 16~mm, $0.791$ against $0.675$ at
25~mm), and the real-label ranking agrees ($0.273$ against $0.232$), but neither
reaches the classical operating point. Denoising is architecture-insensitive:
U-Net, TCN and BiGRU means are $0.364$, $0.365$ and $0.381$. Gap filling
separates them modestly, with the BiGRU best at 5 and 30 days ($2.378$ and
$2.892$~mm) and the TCN second ($2.401$ and $2.979$), but the BiGRU costs
$5.1$~h per seed against $6$~min for the TCN. The seed spread dominates the
architecture difference in detection; in gap filling the ordering is stable
across architectures, which is the stronger form of the result.

\paragraph{Limitations.}
The truth families are still constructed, not observed: the in-family and
out-of-family signals are deterministic functions with known parameters, and the
out-of-family generator is one choice among many. The real-label arm scores
against a catalogue, not against a measured break, and its 3$\sigma$ amplitude
filter removes weak continuity changes that a detector might still be expected
to find. All experiments are single-station; common-mode error and network
filtering are outside the protocol. The downstream arm uses six stations, two
components, three realisations and two U-Net seeds, so its calibration ratios
carry a wide interval and are read as a comparison across methods, not as an
absolute calibration of geodetic uncertainty; the untreated series itself covers
$0.389$ at 1$\sigma$ because of irregular sampling, and that shortfall is shared
by every method. The random-walk arm tests one amplitude ($1.0$~mm\,yr$^{-1/2}$)
and one noise model. The 365-day gap and the 16--25~mm offsets lie at or beyond
the training support and measure extrapolation rather than in-distribution
skill. The external centre arm is descriptive: it scores held-out JPL series
with networks trained on the same generators, not networks retrained at JPL,
so the design is a single primary analysis centre plus JPL, not a multi-centre
evaluation; its JPL group contains 18 component series against 183 for NGL. The
regional split is uneven, with 123 component series in Asia and six in the
residual group, so regional means outside Asia and North America rest on few
stations and are read as a consistency check on the ranking rather than as
regional predictions.
